\section{\sysname Design}
\label{sec:design}

\subsection{Overview}
\label{sec:design:overview}

\sysname turns an upstream activity mask into the coordinate stream of a
tiled convolution. A layer receives a dense feature tensor $\mathbf{x}$ and
its input mask $M_{in}$. It produces the output values $\mathbf{y}$ requested
by the mask policy, the propagated output mask $M$, and a worklist $\mathbf{q}$
of active output coordinates. The operator separates this work into two stages
(Figure~\ref{fig:overview}):

\begin{enumerate}
  \item \textbf{Construction stage.} A mask-space kernel evaluates
        Eq.~\eqref{eq:mask}, writes $M$, and compacts its active coordinates
        into $\mathbf{q}$. It scans the output grid in spatial tiles, so each
        emitted segment retains the local traversal order of one tile.

  \item \textbf{Compute stage.} A tiled convolution kernel reads its output
        coordinates from $\mathbf{q}$ instead of enumerating a contiguous
        output region. For those coordinates, it preserves the regular
        reduction over input channels and kernel offsets used by dense
        convolution.
\end{enumerate}

This separation removes coordinate discovery from the convolution kernel's
critical path and establishes spatial structure before feature data are
accessed. Construction touches only masks and coordinates; compute retains
dense feature and weight tensors and executes one output for every worklist
entry. The following sections describe tiled construction and worklist-driven
compute.

% TODO: overview figure (double column)
\begin{figure*}[t]
  \centering
  \fbox{\parbox{0.9\textwidth}{\centering\vspace{2.5cm}%
    \textbf{[TODO]} \sysname overview. Left: dense input feature tensor and
    input mask. Middle: tiled construction emits the output mask and compact
    worklist segments. Right: worklist coordinates replace the contiguous
    coordinate source of tiled convolution; only active outputs are written.
    \vspace{2.5cm}}}
  \caption{The \sysname two-stage pipeline. Construction discovers active
  output coordinates and compacts each spatial tile into a worklist segment.
  Compute uses the resulting coordinates to drive a tiled convolution while
  retaining its regular channel and kernel reduction.}
  \label{fig:overview}
\end{figure*}

\subsection{Tiled Worklist Construction}
\label{sec:design:construct}

Construction must determine the active output set before convolution without
moving feature data or imposing a device-wide ordering pass. We partition the
output grid into contiguous construction tiles $T_1,\ldots,T_N$, each with a
fixed local traversal. For every position $p\in T_n$, the stage evaluates
Eq.~\eqref{eq:mask} and writes the result to $M[p]$. It then selects the active
positions within that tile:
\begin{equation}
  \mathcal{I}_n = \{j : M[T_n[j]]=1\},
  \label{eq:support}
\end{equation}
and emits their linear output coordinates in local traversal order:
\begin{equation}
  \mathbf{q}_n = (T_n[j])_{j\in\mathcal{I}_n}.
  \label{eq:pertile-gather}
\end{equation}

Tiles execute in parallel and reserve disjoint segments of the final
worklist. If $P$ denotes their reservation order, the resulting array is
\begin{equation}
  \mathbf{q} = [\mathbf{q}_{P_1};\ldots;\mathbf{q}_{P_N}],
  \qquad |\mathbf{q}|=\|M\|_0.
  \label{eq:worklist}
\end{equation}
Parallel scheduling may change the order of these segments, so \sysname does
not require or provide a global raster order. Its guarantee is tile-local:
every $\mathbf{q}_n$ contains coordinates from one contiguous output region
in that tile's traversal order. Generating this order during compaction avoids
a subsequent sort over the active set.

Algorithm~\ref{alg:construct} summarizes the stage. The fused pass visits
$H_{out}W_{out}$ output positions and examines at most $k^2$ input-mask
entries per position. Its $O(H_{out}W_{out}k^2)$ work is independent of
$C_{in}$ and $C_{out}$. We include construction in operator latency and measure
its cost directly in Section~\ref{sec:eval:overhead}.

\begin{algorithm}[t]
\caption{Tiled Worklist Construction}
\label{alg:construct}
\begin{algorithmic}[1]
\REQUIRE Construction tiles $T_1,\ldots,T_N$, input mask $M_{in}$
\ENSURE Worklist $\mathbf{q}$, output mask $M$
\FOR{each tile $T_n$ \textbf{in parallel}}
  \FOR{each position $p\in T_n$}
    \STATE $M[p]\leftarrow$ Eq.~\eqref{eq:mask}
  \ENDFOR
  \STATE $\mathbf{q}_n\leftarrow(T_n[j])_{j:M[T_n[j]]=1}$
  \STATE $b_n\leftarrow\operatorname{reserve}(|\mathbf{q}_n|)$
  \STATE $\mathbf{q}[b_n:b_n+|\mathbf{q}_n|]\leftarrow\mathbf{q}_n$
\ENDFOR
\end{algorithmic}
\end{algorithm}

\subsection{Worklist-Driven Convolution}
\label{sec:design:compute}

Implicit GEMM views convolution as a matrix product whose rows correspond to
output coordinates, columns to output channels, and reduction dimension to
input channels and kernel offsets. A dense kernel derives each row coordinate
from a contiguous output tile. \sysname instead loads those coordinates from
$\mathbf{q}$; the column and reduction dimensions remain regular.

The compute stage partitions the worklist into blocks of $B_M$ coordinates,
the output channels into blocks of $B_N$, and the reduction into steps of
$B_K$. Each compute tile first loads a coordinate block $\mathbf{q}'$. For
every reduction step, it derives input addresses from $\mathbf{q}'$, loads the
corresponding input and weight tiles, and accumulates their products. It then
writes the result to the dense output tensor at the same coordinates
(Algorithm~\ref{alg:compute}). The worklist schedules exactly the support of
$M$.

The input tile in Algorithm~\ref{alg:compute} is loaded directly from
$\mathbf{x}$ into the kernel's working storage; it is not materialized as a
separate global matrix. Construction order makes these loads less irregular
than an arbitrary active-coordinate list. Entries within each
$\mathbf{q}_n$ address nearby output positions and therefore nearby,
frequently overlapping input neighborhoods. A compute tile inherits this
locality when its coordinate block lies within a construction segment. It may
also cross segment boundaries, which affects locality but not correctness.
This limited guarantee lets compute exploit available spatial structure
without requiring globally sorted coordinates.

\begin{algorithm}[t]
\caption{Worklist-Driven Convolution}
\label{alg:compute}
\begin{algorithmic}[1]
\REQUIRE Worklist $\mathbf{q}$, input $\mathbf{x}$, weight $\mathbf{w}$
\ENSURE Output $\mathbf{y}$ at coordinates in $\mathbf{q}$
\FOR{each output tile $(b_m,b_n)$ \textbf{in parallel}}
  \STATE $\mathbf{q}'\leftarrow
    \mathbf{q}[(b_m{-}1)B_M:b_mB_M]$
  \STATE $\mathbf{a}\leftarrow\mathbf{0}_{B_M\times B_N}$
  \FOR{each reduction tile $r$ of width $B_K$}
    \STATE $\mathbf{x}'\leftarrow$ input tile addressed by
      $(\mathbf{q}',r)$
    \STATE $\mathbf{w}'\leftarrow$ weight tile addressed by $(b_n,r)$
    \STATE $\mathbf{a}\leftarrow\mathbf{a}+\mathbf{x}'\mathbf{w}'$
  \ENDFOR
  \STATE $\mathbf{y}[\mathbf{q}',b_n]\leftarrow\mathbf{a}$
\ENDFOR
\end{algorithmic}
\end{algorithm}
